\section*{Chapter \ref{ch:g8:searching-and-sorting} --- Searching and Sorting}

\begin{solution}{exo:g8:searching-and-sorting:1}
Step 1: \code{low} 0, \code{high} 8, \code{middle} 4, \omterm{def:g7:lists:list}{item} 29, smaller than
37, so \code{low} becomes 5. Step 2: \code{low} 5, \code{high} 8,
\code{middle} 6, \omterm{def:g7:lists:list}{item} 44, larger, so \code{high} becomes 5. Step 3:
\code{low} 5, \code{high} 5, \code{middle} 5, \omterm{def:g7:lists:list}{item} 37: found at position 5,
in 3 steps.
\end{solution}

\begin{solution}{exo:g8:searching-and-sorting:2}
Pass 1 swaps the smallest, 1, with 6: \code{[1, 2, 8, 5, 6]}. Pass 2 finds
that the smallest of the rest, 2, is already in place:
\code{[1, 2, 8, 5, 6]}.
\end{solution}

\begin{solution}{exo:g8:searching-and-sorting:3}
Pass 1 slides 2 before 6: \code{[2, 6, 8, 5, 1]}. Pass 2 leaves 8 where it
is, since 6 is smaller: \code{[2, 6, 8, 5, 1]}.
\end{solution}

\begin{solution}{exo:g8:searching-and-sorting:4}
A \omterm{def:g8:searching-and-sorting:linear-search}{linear search} that does not stop at the first match:

\omcode{code/grade-8/04-searching-and-sorting/exercises/all_positions.py}{1}{9}{All
the positions of a \omterm{def:g6:variables-and-input:value}{value}.}\label{lst:g8:searching-and-sorting:sol-all}
\end{solution}

\begin{solution}{exo:g8:searching-and-sorting:5}
\code{low} 0, \code{high} 1, \code{middle} 0: 4 is larger than 2, so
\code{high} becomes 0. Then \code{low} 0, \code{high} 0, \code{middle} 0: 4
is larger again, and \code{high} becomes 0 again. Nothing changes any more,
and \code{low <= high} stays true forever. The middle \omterm{def:g7:lists:list}{item} has already been
compared, so it must be left out: \code{high = middle - 1}. Then each step
makes the part searched smaller, and the \omterm{def:g4:algorithms-repeat-and-decide:loop}{loop} ends.
\end{solution}

\begin{solution}{exo:g8:searching-and-sorting:6}
\code{range(2, len(items))} starts at position 2, so the \omterm{def:g7:lists:list}{item} at position 1
(here 1) is never slid back. Starting at 1, \code{range(1, len(items))},
sorts the \omterm{def:g7:lists:list}{list}: \code{[1, 3, 5]}.
\end{solution}

\begin{solution}{exo:g8:searching-and-sorting:7}
Saved in the \omterm{def:g3:files-and-folders:folder}{folder} of \cref{lst:g8:searching-and-sorting:binary}, the \omterm{def:g4:programs-and-operating-system:program}{program}
imports its \code{binary_search}:

\omcode{code/grade-8/04-searching-and-sorting/exercises/max_guesses.py}{1}{9}{The
most guesses between 1 and 1,000.}\label{lst:g8:searching-and-sorting:sol-max}

It prints \code{at most 10 guesses for 1 to 1000}. Indeed 9 guesses cover
only $2^9 - 1 = 511$ \omterm{def:g7:lists:list}{items}, and 10 cover $2^{10} - 1 = 1{,}023$.
\end{solution}

\begin{solution}{exo:g8:searching-and-sorting:8}
The \omterm{def:g8:searching-and-sorting:selection-sort}{selection sort} of the chapter, comparing \code{len} of the words:

\omcode{code/grade-8/04-searching-and-sorting/exercises/by_length.py}{1}{13}{Words
from the shortest.}\label{lst:g8:searching-and-sorting:sol-length}

It prints \code{['fig', 'kiwi', 'apple', 'banana', 'cherry']}.
\end{solution}

\begin{solution}{exo:g8:searching-and-sorting:9}
Pass 1 compares the smallest found so far with each of the 4 other
\omterm{def:g7:lists:list}{items}; pass 2 with the 3 after the first place, then 2, then 1: $4 + 3 + 2 + 1 = 10$ \omterm{def:g6:if-else:comparison}{comparisons}. It does not depend on the
order: \omterm{def:g8:searching-and-sorting:selection-sort}{selection sort} always looks at the whole rest of the \omterm{def:g7:lists:list}{list} to be sure
it has the smallest.
\end{solution}

\begin{solution}{exo:g8:searching-and-sorting:10}
\code{[5, 1, 3]}: the middle \omterm{def:g7:lists:list}{item} is 1, smaller than 5, so the search goes
right; it finds 3, still smaller, goes right again, and stops with $-1$. \omterm{def:g8:searching-and-sorting:binary-search}{Binary
search} throws away a half without looking at it; that is safe only if the
order guarantees that the target cannot be there.
\end{solution}

\begin{solution}{exo:g8:searching-and-sorting:11}
Pass $i$ (for $i$ from 1 to $n - 1$) slides the \omterm{def:g7:lists:list}{item} at position $i$. On a
\omterm{def:g8:searching-and-sorting:sorted}{sorted list} it is compared once with the \omterm{def:g7:lists:list}{item} before it, which is smaller, and
stays: one \omterm{def:g6:if-else:comparison}{comparison} per pass, $n - 1$ in all. On a reversed \omterm{def:g7:lists:list}{list} it is
smaller than all the $i$ \omterm{def:g7:lists:list}{items} before it, so it is compared with each of them
and slides to the front: $i$ \omterm{def:g6:if-else:comparison}{comparisons}. In all,
$1 + 2 + \dots + (n - 1)$. Writing this sum twice, once forwards and once
backwards, gives $n - 1$ pairs that each add up to $n$, so the sum is
$n(n-1)/2$.
\end{solution}

\begin{solution}{exo:g8:searching-and-sorting:12}
A helper \omterm{def:g7:functions:function}{function} decides which of two books comes first; the \omterm{def:g8:searching-and-sorting:insertion-sort}{insertion sort}
uses it in place of \code{>}:

\omcode{code/grade-8/04-searching-and-sorting/exercises/sort_books.py}{1}{14}{The
books by year, then by title.}\label{lst:g8:searching-and-sorting:sol-books}

A test of the chapter checks it on records with the same year.
\end{solution}

\begin{solution}{pb:g8:searching-and-sorting:1}
\textbf{1.} Guess 8: higher, 9 to 16 left. Guess 12: lower, 9 to 11 left.
Guess 10: higher, 11 left. Guess 11: found, in 4 guesses.

\textbf{2.} 1,000, 500, 250, 125, 63, 32, 16, 8, 4, 2, 1: ten halvings, so
about 10 guesses.

\textbf{3.} Throwing away half of the \omterm{def:g7:lists:list}{list} is safe only if the order says the
target cannot be there.

\textbf{4.} Going right keeps the larger half, since \code{//} rounds down:
it is the case where the most guesses are needed.

\textbf{5.} Each time the size is multiplied by 10, 3 or 4 more guesses are
needed:

\omoutput{code/grade-8/04-searching-and-sorting/outputs/worst_guesses.txt}{The
worst case, from 10 to a million \omterm{def:g7:lists:list}{items}.}\label{out:g8:searching-and-sorting:worst}

Multiplying by 10 is a little more than three doublings, since
$2^3 = 8 < 10 < 16 = 2^4$.

\textbf{6.} 19 guesses cover $2^{19} - 1 = 524{,}287$ \omterm{def:g7:lists:list}{items}, not enough;
20 cover $2^{20} - 1 = 1{,}048{,}575$, which is at least a million.

\textbf{7.} A million: every word is looked at when the word sought is last
or absent.

\textbf{8.} The cost of 1,000 \omterm{def:g8:searching-and-sorting:linear-search}{linear searches} in a mixed \omterm{def:g7:lists:list}{list}, against the
cost of sorting the \omterm{def:g7:lists:list}{list} once and then making 1,000 \omterm{def:g8:searching-and-sorting:binary-search}{binary searches}.

\textbf{9.}

\omoutput{code/grade-8/04-searching-and-sorting/outputs/sort_or_scan.txt}{\omterm{def:g6:if-else:comparison}{Comparisons}
without and with sorting.}\label{out:g8:searching-and-sorting:sort-or-scan}

The \omterm{def:g8:searching-and-sorting:linear-search}{linear searches} cost 500,500 \omterm{def:g6:if-else:comparison}{comparisons}. Sorting costs 499,500 with
\omterm{def:g8:searching-and-sorting:selection-sort}{selection sort} and 252,747 with \omterm{def:g8:searching-and-sorting:insertion-sort}{insertion sort}; the 1,000 \omterm{def:g8:searching-and-sorting:binary-search}{binary searches} add
only 8,987.

\textbf{10.} With \omterm{def:g8:searching-and-sorting:selection-sort}{selection sort}, $499{,}500 + 8{,}987 = 508{,}487$, a little
more than the 500,500 of the \omterm{def:g8:searching-and-sorting:linear-search}{linear searches}: not worth it. With \omterm{def:g8:searching-and-sorting:insertion-sort}{insertion
sort}, $252{,}747 + 8{,}987 = 261{,}734$, about half: worth it.

\textbf{11.} The sort is paid once, and each \omterm{def:g8:searching-and-sorting:binary-search}{binary search} costs about 9
\omterm{def:g6:if-else:comparison}{comparisons} instead of about 500. 100,000 \omterm{def:g8:searching-and-sorting:linear-search}{linear searches} cost about
50 million \omterm{def:g6:if-else:comparison}{comparisons}; 100,000 \omterm{def:g8:searching-and-sorting:binary-search}{binary searches} about 900,000, plus one sort.
Sorting first wins by far.

\textbf{12.} \textbf{Result.} \omterm{def:g8:searching-and-sorting:binary-search}{Binary search} finds any word among 1,000,000
\omterm{def:g8:searching-and-sorting:sorted}{sorted} words in at most 20 \omterm{def:g6:if-else:comparison}{comparisons}, because $2^{20} = 1{,}048{,}576$ is at
least 1,000,000; a \omterm{def:g8:searching-and-sorting:linear-search}{linear search} may need a million. Sorting costs more than
one search, but it is paid once: as soon as the \omterm{def:g7:lists:list}{list} is searched many times,
sorting first wins.
\end{solution}
